\chapter{交互演示与数据发布}
\label{ch:demo}

演示是第 \ref{ch:proof} 章命题的可操作规格，不是另一套模型。积分、更新、评分和投影属于可测试的数值核；界面只读取该核的状态流。页面上的每个量和下载字段都要能指回已经编号的方程。数值核尚未实现时，本章只作为规格。

\section{页面必须同时露出的四块}

\begin{enumerate}
\item \textbf{边界。} 加热、蒸汽、风速、排汽。策略可以写它们，操作者也可以接管。接管的时段标记为干预，路径按 $\mathrm{do}$ 计算，不按策略相关计算。
\item \textbf{状态。} $T_1$、$T_2$、含水、$\alpha_m$、$R_{\mathrm{int}}$ 的后验均值和方差。真值曲线与估计曲线分开着色。
\item \textbf{路径。} 当前蒸汽设定下 A、B、C 三项。符号在抽样中不稳定时只留区间。
\item \textbf{投影。} 品质模块请求的加热量，以及进入边界的加热量。管温贴近 $h_T=0$ 时二者必须分开，否则命题 \ref{prop:shield} 无法在页面上被看见。
\end{enumerate}

策略四档与第 \ref{ch:sim} 章相同：开环、腔温闭环、无探针优化、对偶控制。可改的样品因素是厚度档、初始温度、初始含水和种子。

\section{页面上的机制，而不是说明文字}

默认视图把 P4 画成两段色带，叠在 $T_1$ 与 $R_{\mathrm{int}}$ 方差上。方差不降的运行，停止标记保持在动作前的位置，用来核对命题 \ref{prop:invalid}。另有一个稳态开关：长时间内只维持腔温，页面上 $R_{\mathrm{int}}$ 的方差保持为先验，用来核对命题 \ref{prop:steady}。

厚度拨到信息上限以外时，剩余时间的位置改成区间，不给出单点。这与第 \ref{ch:ati} 章的显示规则是同一条。

\section{下载字段}

一次运行对应三个文件，字段固定，便于换语言重算后逐列对比。

\begin{center}
\begin{tabularx}{\textwidth}{llX}
\toprule
文件 & 行或键 & 必须含有 \\
\midrule
\texttt{trajectory.csv} & 每个时间步一行 & $t,T_1,T_2,\hat T_2,m,\alpha_m,R_{\mathrm{int}},\mathrm{Var},u,u_{\mathrm{proj}},\mathrm{mode}$ \\
\texttt{summary.json} & 一次运行 & 停止时刻、$e_{\alpha}$、能量、A、B、C、探针是否有效 \\
\texttt{parameters.json} & 一次运行 & \eqref{eq:twonode} 的全部常数、噪声、种子、策略名、是否同构真值 \\
\bottomrule
\end{tabularx}
\end{center}

\texttt{mode} 取值为品质、信息、保温、干预。参数文件里的数就是积分用的数，不另存展示用副本。

\section{计算与界面分开}

\begin{center}
\begin{tabularx}{\textwidth}{lXX}
\toprule
层 & 职责 & 不允许的越界 \\
\midrule
数值核 & 两节点真值、噪声、集合更新、探针评分、路径三项、投影 & 为了画面平滑而改方程 \\
断言 & 命题 \ref{prop:steady}、\ref{prop:pulse}、\ref{prop:budget}、\ref{prop:shield} 的可重复检查 & 依赖界面点击才成立 \\
界面 & 四块区域、色带、区间、下载按钮 & 自行积分或另存一套参数 \\
离线真值 & 多孔介质及 GAB、各向异性、脂肪夹层、收缩 & 改列名后再进界面 \\
\bottomrule
\end{tabularx}
\end{center}

默认真值与在线模型同构，只承担命题核对。多孔介质轨迹沿用同一列名追加。同构结果不标成主结果。界面不展示 \texttt{trajectory.csv} 中不存在的百分比。主对照仍以第 \ref{ch:sim} 章的日志为准。

演示同时包含机制与仿真，二者使用同一时刻的状态，不分成两个互不引用的页面。

\begin{itemize}
\item 数值核用 Python、NumPy 与 SciPy。它输出第 \ref{ch:demo} 章的字段，并在无界面模式下完成命题断言。
\item 机制视图用矢量图显示父节点、路径 A/B/C、探针两段和投影。线的明暗由当前偏效应或方差决定，不使用与状态无关的装饰动画。
\item 仿真视图用三维场景。食品几何按当前厚度与收缩显示，切面颜色表示温度或进度，切面位置可由操作者移动。腔体只显示加热、蒸汽、风和排汽四个边界的状态。
\item 三维视图读取数值核的场，不在着色器里另做传热。两节点阶段，切面由表面节点与内部节点插值并标明“两节点插值”。多孔介质日志接入后，切面改用该日志的空间列。
\item 界面用 React 与 TypeScript。三维渲染使用 WebGL 组件，只消费数值核已经算出的数组。
\item 本地接口传递字段。下载文件仍由数值核序列化。
\end{itemize}

无界面测试先于三维视图。视图接入后，用同一组种子比对：界面导出的 CSV 与数值核直接写出的 CSV 逐列一致。不一致则演示不合格，即使画面完整。

\section{三维场景里必须能读出的量}

场景不是腔体效果图。每个可见元素对应一个状态分量。

\begin{itemize}
\item 食品实体按 $L(t)$ 缩放。切面滑块给出厚度坐标上的 $T$、$\varepsilon_l$ 或 $\alpha$，三者可切换，色标固定在该次运行的最小与最大，不随帧重归一化到全色。
\item 脂肪夹层在 T2 日志中存在时显示为独立薄层；T1 没有该层时不绘制。
\item 蒸汽、排汽、风速、加热用四个边界标记的强度表示，强度等于投影后的 $u$，不是请求值。
\item 探针两段在时间轴上着色，并在食品表面闪动一次；闪动不改变几何。
\item 路径 A/B/C 放在切面旁边的固定图，条的长度是偏效应。符号不稳定时条变为中性色并标区间。
\item 请求加热与投影加热分成两根针。两针重合时表示屏障未动作。
\item 相机只提供旋转、缩放和切面。没有与状态无关的漫游动画。
\end{itemize}

演示数据预置 S1 的一个种子、S2 的一次遮挡、S2 的一次屏障贴近，以及 T1 与 T2 各一条同名种子。操作者可改策略并重算两节点；多孔介质条目从日志载入，不在浏览器内重解 PDE。
